\section{Methods}
\label{sec:methods}

We propose a method to systematically uncover Simpson's paradox for trends in data. We denote as $Y$ the dependent variable in the data set, i.e., an outcome being measured, and as $\mathbf{X} = \{X_1, X_2,\dots,X_m\}$ the set of $m$ independent variables or features. The goal of the method is to identify pairs of variables $(X_p, X_c)$---\emph{Simpson's pairs}---such that %$Y$ shows an upward (resp. downward) trend as a function of $X_p$, but
a trend in $Y$ as a function of $X_p$ disappears or reverses when the data is disaggregated by conditioning on $X_c$. %then $Y$ conversely shows a downward (resp. upward) trend as a function of $X_p$.
More specifically, our method searches for pairs of variables $(X_p, X_c)$ such that
\begin{align}
	&\frac{d}{dx_p}\mathbb{E}[Y|X_p = x_p] > 0\;\;\;\; \forall x_p, \label{eq:trend_sp1} \\
	&\frac{d}{dx_p}\mathbb{E}[Y|X_p = x_p, X_c = x_c] \leq 0 \;\;\;\; \forall x_p, x_c.
	\label{eq:trend_sp2}
\end{align}
and vice versa (i.e., $d\mathbb{E}[Y|X_p = x_p]/dx_p < 0$, $d\mathbb{E}[Y|X_p = x_p, X_c = x_c]/dx_p \geq 0$). Equations~\eqref{eq:trend_sp1} and \eqref{eq:trend_sp2} hold if the expected value of $Y$ is a monotonically increasing (or decreasing) function of $X_p$ alone, but conditioned on $X_c$ is a monotonically decreasing (resp. increasing) function of $X_p$, or is constant.

\subsection{Finding Simpson's Pairs}
With this goal in mind, we employ linear models to quantify the relationship between $Y$, an independent variable $X_p$, and a conditioning variable $X_c$ upon which the data is disaggregated. Firstly, on the aggregate level, we model the relationship between $Y$ and $X_p$ as a linear model of the form
\begin{equation}
	\mathbb{E}[Y|X_p = x_p] = f_p(\alpha + \beta x_p),
	\label{eq:f_p}
\end{equation}
where $f_p(\alpha + \beta x_p)$ is a monotonically increasing function of its argument $\alpha + \beta x_p$. The parameter $\alpha$ in Eq.~\eqref{eq:f_p} is the intercept of the regression function, while the trend parameter $\beta $ quantifies the effect of $X_p$ on $Y$. Secondly, for the disaggregation, we fit linear models of the form of Eq.~\eqref{eq:f_p} but with different values of the parameters $\alpha$ and $\beta $ depending on the value of $X_c$:	
\begin{equation}
	\mathbb{E}[Y|X_p = x_p, X_c = x_c] = f_{p,c}(\alpha(x_c) + \beta (x_c)x_p),
	\label{eq:f_pc}
\end{equation}

%Although it is possible to fit a simpler linear model of both $X_p$ and $X_c$ combined as $f_{p,c}(\alpha + \beta X_p + \beta_cX_c)$, this can be limiting because of the assumption that the relationship between $X_p$ and $Y$ is the same (as quantified by $\beta $) for all values of $X_c$. \hl{For instance, for $X_c = a$, the formula become $f_{p, c}(\alpha + a \times \beta_c + \beta X_p)$, which shows that for all subgroups we have the same trend coefficients $\beta$ and have intercepts that vary linearly with $a$. In our proposed method, however, each group can have different intercept and coefficient, which makes finding paradox pairs in heterogeneous data more flexible (we will show an example of the outputs of the two methods in Results section)}. Here, we bin the values of $X_c$ and in each bin we fit a linear model of the form of Eq.~\eqref{eq:f_pc}. This approach leads to unique parameters $(\alpha(x_c), \beta (x_c))$ for each bin of $X_c$, quantifying the relationship between $Y$ and $X_p$ conditioned on the value of $X_c$. Indeed this latter approach is necessary - from our results (Figs.~\ref{fig:fig2} and \ref{fig:paradox2b}) it is clear that the trend parameters $\beta (x_c)$ of the fitted lines vary significantly depending on the value $x_c$.

When fitting linear models $f(\alpha + \beta X)$ we have not only a fitted trend parameter $\beta$ but also a $p$-value which gives the probability of finding an intercept $\beta$ at least as extreme as the fitted value under the null hypothesis $H_0:\beta = 0$. From this, we have three possibilities:
\begin{itemize}
  \item $\beta$ is not statistically different from zero (sgn$(\beta)$ = 0), % that
  \item $\beta$ is statistically different from zero and positive (sgn$(\beta)$ = 1), % or that
  \item $\beta$ is statistically different from zero and negative (sgn$(\beta)$ = -1).
\end{itemize}
This mechanism allows us to test for Simpson's paradox by comparing the sign of $\beta$ from the aggregated fit (Eq.~\eqref{eq:f_p}) with the signs of the $\beta$'s from the disaggregated fits (Eq.~\eqref{eq:f_pc}). %\hl
{Although} Eqs.~\eqref{eq:f_p} and \eqref{eq:f_pc} %\hl
{state that the signs from the disaggregated curves should all be different from the aggregated curve, in practice this is too strict, especially as human behavioral data is noisy. Thus, we compare the sign of the fit to aggregated data to the simple average of the signs of fits to disaggregated data--- if the signs are different then we have uncovered an instance of Simpson's paradox}. The summary of the algorithm is the following:

% Moved to results
%In Section~\ref{sec:results} we apply our methodology to data from the Stack Exchange online forum. Here, our dependent variable $Y$ is binary, denoting whether or not an specific answer to a question was accepted as the best answer. In this case of binary outcomes we use the logistic regression linear model of the form
%\begin{equation}
%	f(\alpha + \beta x) = \frac{1}{1+e^{-(\alpha + \beta x)}}.
%\end{equation}
%The parameters $\alpha$ and $\beta$ are fitted using Maximum likelihood, while test of the null hypothesis $H_0:\beta = 0$ is performed using the Likelihood Ratio Test~\cite{casella2002statistical}.

\lstset{language=Python,
  showspaces=false,
  showtabs=false,
  breaklines=true,
  numbers=left,
  stepnumber=1,
  frame=single,
  showstringspaces=false,
  breakatwhitespace=true,
  escapeinside={(*@}{@*)},
  commentstyle=\color{greencomments},
%  keywordstyle=\color{bluekeywords},
  stringstyle=\color{redstrings},
  basicstyle=\ttfamily
}

\begin{lstlisting}[title=Trend Simpson's Paradox Algorithm][language=Python]
def trend_simpsons_pair(X, Y):
 paradox_pairs = []
 for paradox_var in vars:
  beta, pvalue = trend_analysis(X[paradox_var], Y)
  agg = sgn(beta, pvalue)
  for condition_var in vars:
   if paradox_var != condition_var:
    dagg = []
    for con_gr in bins_of(condition_var):
     beta, pvalue = trend_analysis(X[paradox_var | con_gr], Y)
     dagg.append(sgn(beta, pvalue))
    if agg != sgn(mean(dagg)):
     paradox_pairs.append([paradox_var, condition_var])
 return paradox_pairs

 def sgn(beta, pvalue = 0.0):
  return (0 if (beta == 0 or pvalue > 0.05) else (1 if beta > 0 else -1))
\end{lstlisting}

\paragraph{Data Disaggregation}
A critical step in our method is disaggregating data by conditioning on variable $X_c$. The idea behind disaggregation is to segment data into more homogeneous subgroups of similar elements. For multinomial variables $X_c$, disaggregation step simply involves grouping data by unique values of $X_c$. However, for continuous $X_c$ or discrete variables with large range, this step is more complex. We can bin the elements according to their values of $X_c$, but the decision has to be made how large each bin is, whether bin sizes scale linearly or logarithmically, etc. If the bin is too small, it may not contain enough samples for a statistically significant measurement, but if it is too large, the samples may be too heterogeneous for a robust trend. %hl
{In our experiments described below, bins of fixed size successfully identify Simpson's pairs, though we realize that more sophisticated binning techniques can allow us to isolate more pairs or reduce the number of false positives.}

\subsection{Mathematical Analysis of the Paradox}
\label{sec:math}
We have presented a mathematical formulation of Simpson's paradox in terms of the derivatives of conditional expectations as given by Eqs.~\eqref{eq:trend_sp1} and \eqref{eq:trend_sp2}, and
we now examine these equations to get a better insight into the origins and causes of this paradox.

The expectation in Eq.~\eqref{eq:trend_sp1} can be related to that of Eq.~\eqref{eq:trend_sp2} as
\begin{eqnarray}
% \nonumber % Remove numbering (before each equation)
	& & \mathbb{E}[Y|X_p =  x_p] = \label{eq:cond_exp} \\ \nonumber
 & & \int_{X_c}\mathbb{E}[Y|X_p = x_p, X_c = x_c]\Pr(X_c=x_c|X_p=x_p)dx_c,
\end{eqnarray}
and differentiating this expectation w.r.t. $x_p$ allows us to compare the trends of Eqs.~\eqref{eq:trend_sp1} and \eqref{eq:trend_sp2}. The derivative of the right hand side of Eq.~\eqref{eq:cond_exp} with respect to $x_p$ is
\begin{align}
	\int_{X_c}\left(\frac{d}{dx_p}\mathbb{E}[Y|X_p = x_p, X_c = x_c]\right)\Pr(X_c=x_c|X_p=x_p)dx_c + \nonumber \\
	\int_{X_c}\mathbb{E}[Y|X_p = x_p, X_c = x_c]\left(\frac{d}{dx_p}\Pr(X_c=x_c|X_p=x_p)\right)dx_c.
\label{eq:exp_der2}
\end{align}
If $\mathbb{E}[Y|X_p = x_p, X_c = x_c]$ is a non-increasing function of $x_p$---as in Eq.~\eqref{eq:trend_sp2}---then the first integral in Eq.~\eqref{eq:exp_der2} will be non-positive. Thus for $\mathbb{E}[Y|X_p = x_p]$ to be an increasing function of $x_p$, i.e., for Eq.~\eqref{eq:exp_der2} to be positive, the second integral must be positive.

This inequality condition leads to two necessary conditions for the occurrence of Simpson's paradox. The first condition is that
\begin{equation}
	\frac{d}{dx_p}\Pr(X_c=x_c|X_p=x_p) \neq 0,
        \label{eq:SPcond1}
\end{equation}
i.e., the distribution of the conditioning variable $X_c$ is not independent of $X_p$ and so the two variables are correlated. As $X_p$ changes, the distribution of the values of $X_c$ must also change. In the case that the distribution of $X_c$ is independent of $X_p$, then $d\Pr(X_c=x_c|X_p=x_p)/dx = 0$ and so the second integral of Eq.~\eqref{eq:exp_der2} will be zero resulting in no Simpson's paradox.
%This dependency can be clearly seen in the joint distributions of the Stack Exchange paradox pair ($X_p$=\emph{Number of Answers}, $X_c$=\emph{Reputation}) in Fig.~\ref{fig:fig3b}---as $X_p$ increases then the distribution of $X_c$ shifts to increasing values.

The second necessary condition for the occurrence of Simpson's paradox is that the expectation of $Y$, conditioned on $X_p$, must not be independent of $X_c$, i.e.,
\begin{equation}
	\mathbb{E}[Y|X_p = x_p, X_c = x_c] \neq \mathbb{E}[Y|X_p = x_p].
        \label{eq:SPcond2}
\end{equation}
For any given value of $X_p$, the expectation of $Y$ must vary as a function of $X_c$. If the condition of Eq.~\eqref{eq:SPcond2} is not met then the second integral in Eq.~\eqref{eq:exp_der2} becomes
\begin{align}
\label{eq:SP_cond2b}
	&\int_{X_c}\mathbb{E}[Y|X_p = x_p]\left(\frac{d}{dx_p}\Pr(X_c=x_c|X_p=x_p)\right)dx_c \\ \nonumber
	&= \mathbb{E}[Y|X_p = x_p]\frac{d}{dx_p}\left(\int_{X_c}\Pr(X_c=x_c|X_p=x_p)dx_c\right) = 0,
\end{align}
and so Simpson's paradox will not occur. %Again, this condition can be seen to be satisfied in both of the Stack Exchange paradox pairs (Figs.~\ref{fig:fig2} and \ref{fig:paradox2b}) as for any given value of $X_p$ the acceptance probability increases as a function of $X_c$.

Thus, this mathematical analysis has given us an insight into causes for Simpson's paradox in data ---correlations between independent variables and the fact that the distribution of the conditioning variable $X_c$ changes at a faster rate with respect to the independent paradox variable $X_p$ than does the expectation of $Y$. This point will be covered in greater detail in the next section.

%Consider, for example, the (\emph{answer position}, \emph{session length}) paradox pair whose behaviour is illustrated in Fig.~\ref{fig:session}. Although the acceptance probability is decreasing as a function of \emph{answer position} for each value of \emph{session length}, the probability mass of \emph{session length} is constantly moving towards larger values as \emph{answer position} increases. Notice that as \emph{answer position} increments from $a$ to $a+1$, sessions of length $a$ no longer count (as the min session length is now $a+1$). Thus, while \emph{session length} has probability mass $\Pr(X_c=a|X_p=a)$ when $X_p=a$, it has probability mass $\Pr(X_c=a|X_p=a+1) = 0$ at $X_p=a+1$ and so
%\begin{equation}
%  \frac{d}{dx_p}\Pr(X_c=a|X_p=x_p)\|_{x_p=a} = -\Pr(X_c=a|X_p=a).
%\end{equation}
%Meanwhile, for all other values of $X_c$ greater than $a$, the probability mass at $X_p=a+1$ is the same as that at $X_p=a$ (as the number of datapoints is constant along session lenghts) but rescaled to account for the sessions of length $a$ , i.e.,
%\begin{equation}
% \Pr(X_c=x_c|X_p=a+1) = \frac{\Pr(X_c=x_c|X_p=a)}{1-\Pr(X_c=a|X_p=a)}.
%\end{equation}
%and so the rate of change of these probability masses with respect to $X_p$ is
%\begin{eqnarray}
%  &&\frac{d}{dx_p}\Pr(X_c=x_c|X_p=x_p)\|_{x_p=a}  =   \\\nonumber &&\left(\frac{1}{1-\Pr(X_c=a|X_p=a)}-1\right)\Pr(X_c=x_c|X_p=a).
%\end{eqnarray}
%The probability mass function $\Pr(X_c=x_c|X_p=a)$ decreases for the $X_c=a$ corresponding to the smallest value of acceptance probability, while increasing for all other values $X_c=a$. Moreover, the rate of increase of this probability mass is greater than the rate at which the acceptance probability is decreasing, and this causes an upward trend to appear when the data is aggregated.


%\subsection{General Idea}
%Assume we have a independent variables $X$, and dependent variable $Y$, and we want to analyze trends of $Y$ in terms of variables of $X$.
%Our method wants to find two variables, named $p$ and $c$, which the dependent variable $Y$ in terms of $X_p$ has uptrend(/downtrend), but when the data is disaggregated into different subgroups based on value of variable $X_c$, then $Y$ has a downtrend(/uptrend) in terms of $X_p$. This is the instance of Trend-Simpson's paradox.
%The following algorithm can be used to find these type of pairs in variables of a data set. It considered every possible pairs of variable in the dataset as \textit{Paradox Variable}, $p$ and \textit{Conditioning Variable}, $c$, and check whether with these two variables Trend-Simpson's Paradox exists or not.

%
%\lstset{language=Python,
%    basicstyle=\footnotesize,
%    showspaces=false,
%    numbers=left,
%}
%
%\renewcommand{\lstlistingname}{Simpson's Trend Paradox Algorithm}
%\begin{lstlisting}[frame=single][caption=Simpson's Trend Paradox Algorithm]
%Input: X and Y
%
%paradox_pairs = []
%for paradox_var in vars:
% sign_agg = trend_analysis(X[paradox_var], Y)
% for condition_var in vars:
%  if paradox_var != condition_var:
%   sign_disagg = trend_analysis(X[paradox_var | condition_var], Y)
%   if sign_agg != sign_disagg:
%    good_pairs.append([paradox_var, condition_var])
%
%Output: good_pairs
%\end{lstlisting}
%
%In this pseudo code, \textit{Trend Analysis} function could be any algorithm for trend analysis, it just have to return a sign of trend(is it uptrend or downtrend or no-trend?) and goodness of fit for that trend analysis algorithm(Does trend analysis algorithm fits to the data statistically significant or not?). The return value of the algorithm, is all pairs of variables in data set, which is the instance of Trend-Simpson's paradox.
%
%
%
%\subsection{Trend Analysis Algorithm}
%Several different algorithms like ANOVA, Mann-Kendall trend test, Regression analysis are used for trend analysis. Selecting the best algorithm for finding trends in the data depends on the type of dependent variable(binary categorical or continues) and types of trends of dependent variable in terms of independent variables(linear, quadric or cubic).
%
%In Stack Exchange data, we want to study the \textit{Acceptance Probability} of answers. Thus, the dependent variable is a binary variable, \textit{accepted}, which shows whether answer to the question is accepted or not. Logistic regression is appropriate for binary dependent variables. Thus, we use Logistic Regression for trend analysis. In our algorithm, we only consider one variable $X_p$ as dependent variable, so Logistic Regression has the form
%\begin{equation} \label{eq:1}
%F(x) = \frac{1}{1 + e^{-(\beta_{0} + \beta_{1} x)}}
%\end{equation}
%
%$\beta_{0}$ is the intercept and $\beta_{1}$ is coefficient for variable $X_p$ of Logistic Regression. In equation \ref{eq:1}, by increasing of $x$, value of $F(x)$ increase(/decreases) if the sign of $\beta_{1}$ is positive(/negative). So, the positive(/negative) sign of $\beta_{1}$ describes the uptrend(/downtrend) of $Y$ in terms of $X_p$.
%
%%http://onlinelibrary.wiley.com/store/10.1002/sim.2023/asset/2023_ftp.pdf?v=1&t=j5irdlhp&s=48a4ae2316ebfd4eb55b1e7c5bb304767e25e84a
%%http://krex.k-state.edu/dspace/bitstream/handle/2097/530/YingLiu2007.pdf?sequence=1&isAllowed=y
%%http://www2.stat.duke.edu/~zo2/dropbox/goflogistic.pdf
%
%\subsection{Goodness of fit}
%Measuring goodness of fit for Logistic Regression, is a controversial topic. Several articles compare different goodness of fit measures for Logistic regression. \cite{}
%
%Pearson's chi-square $\chi^{2}$ test, Hosmer and Lemeshow's $\hat{C}$ and $\hat{H}$ tests and Deviance D test, are tests used to measure goodness of fit of Logistic Regression.
%Some of these tests compare the overall difference between the observed and fitted values. Among these tests, Pearson's chi-square and Deviance D test are used most often. \cite{} These tests are based on covariate patterns, a unique combination of values of independent variables.
%
%We use $\chi^{2}$ as a measure for goodness of fit for Logistic Regression. Even though, the $\chi^{2}$ is sensitive to sample size. Based on \cite{}, the goodness of fit better approximated by chi-square distribution when the observed number of covariate patterns approximately equals the sample size. To come up with huge number of data points in each group, we use sampling. For each group, we consider 100 data points and keep the ratio between accepted / not accepted answers in each group.
